\section{\bt{Conclusion}{结论}}
\label{sec:limitations}
\bi{We presented \system{}, a shared layer over relational dependency graph $\mathcal G=(\mathcal M,\mathcal H)$ for agents and PostgreSQL: \emph{graph-based verifiable relational kernel extraction} composes accepted sessions over verified join primitives into kernel-$I$ definitions (C1); \emph{lightweight adaptive semantic resolution} enables LLM-scan-free set-valued \emph{Lookup} (C2); and \emph{snapshot-validated definition repair} with \emph{snapshot-bound execution} maintains $\mathcal G$ under P1--P3 via version-coalesced checks, history-preserving repair, and single-snapshot queries (C3). Across TPC-DS, RelBench, and four update benchmarks, \system{} cuts reuse tokens \ReuseTokMinX--\ReuseTokMaxX$\times$ and latency \ReuseLatMinX--\ReuseLatMaxX$\times$ vs.\ \emph{No memory}, hits \LkDirectPrec\% direct lookup precision, matches unseen bindings, and serves zero wrong definition-served answers.}
{本文提出面向智能体与 PostgreSQL 的共享记忆层 \system{}，在关系依赖图 $\mathcal G=(\mathcal M,\mathcal H)$ 上实现：\emph{基于图的可验证关系核提取技术}依托连接原语将采纳会话组装为关系核 $I$（C1）；\emph{轻量级自适应语义解析技术}实现免大模型扫目录的集合查找（C2）；\emph{基于快照验证的定义修复技术}与\emph{快照绑定执行}在 P1--P3 下通过按版本合并验证、历史保持单编辑修复与单快照执行维持有效性（C3）。在 TPC-DS、RelBench 与四类更新基准上，\system{} 较无记忆分别降低复用 token \ReuseTokMinX--\ReuseTokMaxX{} 倍与延迟 \ReuseLatMinX--\ReuseLatMaxX{} 倍，直接命中率 \LkDirectPrec\%，未见绑定全对且更新下定义服务零错答。}

\bi{Three limits remain: (1)~lookup tests a 26-definition catalog and dips to 88\% on unseen-family rewrites, inviting multi-family synthesis and 100-definition organic logs; (2)~conditions miss untouched-history premise shifts under P1--P3, needing lightweight $B$ updates; and (3)~kernel DAGs handle single-fact many-to-one subgraphs rather than multi-fact graphs.}
{三项边界指引后续工作：（1）检索聚焦 26 条目录且跨家族改写降至 88\%，需拓展多家族合成与百级真实日志；（2）数据检查不捕获维持未变历史的 P1--P3 纯语义偏移，需轻量化更新 $B$；（3）将关系核由单事实多对一有向无环图扩展至多事实图。}
